The project implements network slicing in a virtualised SDN environment deployed on an \textbf{Ubuntu 22.04 LTS virtual machine}.  
The experimental setup is based on:
\begin{itemize}
    \item \textbf{Mininet}, used to emulate hosts, switches, and links;
    \item \textbf{Open vSwitch} (via Mininet), used as the software data plane;
    \item \textbf{Python 3} scripts, used to define topologies and automate tests;
    \item \textbf{Ryu controller}, used to implement SDN control logic and OpenFlow rules.
\end{itemize}

The implementation is organised into the following sections:
\begin{enumerate}
    \item \textbf{Topology-slicing}: restriction of communication paths between specific hosts;
    \item \textbf{Service-slicing}: differentiated treatment of selected traffic classes.
\end{enumerate}

\noindent
For each section, validation tests are performed (e.g., connectivity and traffic checks) to evaluate the effectiveness of the proposed slicing policies.